%------------------------------------------------------------------------------%
\section{Funções}

As funções surgem quando uma quantidade (Grandeza) depende de outra. Por
exemplo, podemos dizer que a área de um quadrado está relacionada ao
comprimento do lado, isso quer dizer que uma vez conhecido o comprimento do
lado, $x$, conseguimos calcular a área usando a função $f$ descrita pela
relação $f(x)=x^2$.

É possível representar uma função de quatro maneiras:
\begin{description}
  \item[verbalmente]     descrevendo-a com palavras
  \item[numericamente]   por meio de uma tabela de valores
  \item[visualmente]     através de um gráfico
  \item[algebricamente]  utilizando-se uma fórmula explícita
\end{description}

%------------------------------------------------------------------------------%
\begin{definicao}{Função}{def:funcao}
  Uma \emph{função}\index{Função} $f\colon D \to C$
  é uma relação que associa a cada elemento $x$
  de um conjunto $D$, chamado \emph{domínio}\index{Domínio},
  um \emph{único} elemento
  $f(x)$ (ou $y$) de um conjunto $C$, denominado
  \emph{contradomínio}\index{Contradomínio}.
\end{definicao}
%------------------------------------------------------------------------------%

%------------------------------------------------------------------------------%
\begin{observation}
  Quando o domínio de uma função $f$ não for especificado, então o
  \emph{domínio} (natural)\index{Domínio!natural}
  de $f$ é o conjunto de todos os valores de $x$
  para os quais $f$ está definida.
\end{observation}
%------------------------------------------------------------------------------%

%------------------------------------------------------------------------------%
\begin{observation}
  Devemos ter cuidado ao analisarmos o domínio de uma função, quando essa é
  originada de um problema real.
  Se, por exemplo, a função $f(x)=x+1$, $f(x)$
  representa o valor a ser pago em uma corrida de táxi em função do número de
  quilômetros rodados $x$, então o domínio de $f$ será
  \[
    D= \left\{ x \in \R \st x \geq 0 \right\},
  \]
  uma vez que o número de quilômetros rodados é sempre não-negativo.
\end{observation}
%------------------------------------------------------------------------------%

Note que o contradomínio $C$ de uma função $f$ pode ser qualquer conjunto que
contenha, dentre outros elementos, os valores de $f(x)$. Sendo assim, em muitas
situações estamos interessados em trabalhar apenas com o conjunto que contém
apenas os valores gerados por $f$, ao qual damos o nome de
\emph{conjunto imagem}.

%------------------------------------------------------------------------------%
\begin{definicao}{Imagem}{def:imagem}
  Dada uma \emph{função} $f\colon D\to C$, o valor $f(x)$ para $x\in D$ é
  chamado \emph{imagem}\index{Imagem} de $x$.
  O conjunto de todas as imagens da função $f$ é
  denominado \emph{conjunto imagem}\index{Conjunto!Imagem}
  (ou simplesmente imagem) e denotado por
  $\Im(f)$, ou seja,
  \[
    \Im(f) = \left\{ f(x) \st x \in D \right\}.
  \]
\end{definicao}
%------------------------------------------------------------------------------%

%------------------------------------------------------------------------------%
\begin{definicao}{Gráfico}{def:grafico}
  O \emph{gráfico}\index{Gráfico} de $f$ é o conjunto dos pontos $(x, y)$
  do plano cartesiano
  que satisfazem $y=f(x)$, ou seja,
  \[
    G(f) = \left\{ \big(x, f(x)\big) \st x\in D \right\}.
  \]
\end{definicao}
%------------------------------------------------------------------------------%

%------------------------------------------------------------------------------%
\begin{figure}
  \centering
  \includegraphics[scale=0.5]{funcao1}
  \caption{Figura retirada do livro Cálculo 1, James Stewart.}
\end{figure}
%------------------------------------------------------------------------------%

%------------------------------------------------------------------------------%
\subsection{Quando uma aplicação não é função?}

%------------------------------------------------------------------------------%
\begin{example}
  Verifique se a equação $y^2+x^2-16=0$ permite uma definição
  de $y$ como função de $x$.

  \solution
  Inicialmente vamos ver se conseguimos explicitar $y$ como função de $x$.
  \begin{align*}
    y^2 & = 16-x^2 \\
    y   & = \pm \sqrt{16-x^2}
  \end{align*}
  Logo, para cada $x \in \R$, temos dois valores para $y$. Portanto, a
  equação $y^2+x^2-16=0$ não permite uma definição de $y$ como função de $x$.
\end{example}
%------------------------------------------------------------------------------%

Nos exemplos a seguir vamos explorar como o domínio natural das funções
pode ser estudado.

%------------------------------------------------------------------------------%
\begin{example}
  Faça o estudo do domínio da Função
  \(
    f(x) = 1-x^4
  \)

  \solution
  Note que $f$ está bem definida para todo $x \in \R$, logo
  \[
    D(f) = \R.
  \]
\end{example}
%------------------------------------------------------------------------------%

%------------------------------------------------------------------------------%
\begin{example}
  Faça o estudo do domínio da Função
  \(
    g(x) = \frac{1}{x+4}
  \)

  \solution
  Nesse caso temos que $g$ está definida quando
  \begin{align*}
    x+4 & \neq 0 \\
    x   & \neq -4,
  \end{align*}
  logo
  \[
    D(g) = \R-\left\{-4\right\}
  \]
  ou
  \[
    D(g) = \left\{ x \in \R \st x \neq -4 \right\}.
  \]
\end{example}
%------------------------------------------------------------------------------%

%------------------------------------------------------------------------------%
\begin{example}
  Faça o estudo do domínio da Função
  \(
    h(x) = \sqrt{9-x}
  \)

  \solution
  Para que $h$ seja bem definida
  \begin{align*}
    9-x & \geq 0 \\
    x   & \leq 9,
  \end{align*}
  logo
  \[
    D(h) = \left\{ x \in \R \st x \leq 9 \right\}.
  \]
\end{example}
%------------------------------------------------------------------------------%

%------------------------------------------------------------------------------%
\begin{example}
  Faça o estudo do domínio da Função
  \(
    f(t) = \frac{\sqrt{2+t}}{\sqrt{2t-6}}
  \)

  \solution
  Para que $f$ esteja bem definida teremos que ter
  \begin{align*}
    2+t & \geq 0 \\
    t   & \geq -2
  \end{align*}
  e
  \begin{align*}
    2t-6 & > 0 \\
    t    & > 3,
  \end{align*}
  fazendo a interseção desses intervalos temos que
  \[
    D(f) = \left\{ x \in \R \st x \geq 3 \right\}.
  \]
\end{example}
%------------------------------------------------------------------------------%

%------------------------------------------------------------------------------%
\begin{example}
  Faça o estudo do domínio da Função
  \(
    f(x) = \sqrt[3]{2x-3}
  \)

  \solution
  Como o índice a raiz é ímpar, então
  \[
    D(f) = \R.
  \]
\end{example}
%------------------------------------------------------------------------------%

%------------------------------------------------------------------------------%
\begin{example}
  Faça o estudo do domínio da Função
  \(
    p(x) = \frac{\sqrt{x-1}}{x^2-5x+6'}
  \)

  \solution
  Para que $p$ esteja bem definida teremos que ter
  \begin{align*}
    x-1 & \geq 0 \\
    x   & \geq 0
  \end{align*}
  e
  \[
    x^2-5x+6 \neq 0
  \]
  que corresponde a
  \[
    x \neq 2 \quad\text{ e }\quad
    x \neq 3,
  \]
  portanto
  \[
    D(p) = \left\{ x \in \R \st x \geq 1,\; x \neq 2, x \neq 3 \right\}.
  \]
\end{example}
%------------------------------------------------------------------------------%
